\section{RUANG LINGKUP DAN BATASAN PERANGKAT LUNAK}

Bagian ini mendefinisikan ruang lingkup \textit{Minimum Viable Product} (MVP) Savora, yaitu
batas fungsionalitas yang menjadi fokus pengembangan untuk lomba CODE 6.0. Ruang lingkup ini
mengacu pada dokumen kebutuhan produk (PRD) Savora dan diposisikan sebagai rancangan MVP yang
dapat diprototipekan dan diuji, bukan klaim bisnis yang sudah berjalan penuh.

\subsection{Ruang Lingkup MVP (In-Scope)}

Fungsionalitas berikut termasuk dalam ruang lingkup MVP Savora. Untuk transparansi, status
implementasi ditandai sebagai \emph{sudah berjalan} (telah ada pada basis kode) atau
\emph{direncanakan} (menjadi bagian MVP dan sedang/akan diimplementasikan).

\begin{enumerate}[leftmargin=*]
    \item \textbf{Autentikasi dan Manajemen Role} : Registrasi dan login untuk empat peran ---
    Customer, UMKM, Admin, dan Mitra Donasi --- dengan \textit{role-based access control}.

    \item \textbf{Manajemen Produk UMKM} : UMKM menambah, mengubah, dan menghapus produk makanan
    surplus (produk reguler maupun \textit{mystery box}) lengkap dengan harga asli, harga rescue,
    stok, foto, alamat pickup, dan batas waktu (\textit{expires\_at}).

    \item \textbf{Food Trust Index} : Penilaian kelayakan berbasis aturan dari input UMKM
    (kategori, jam produksi, metode penyimpanan, kondisi kemasan) yang menghasilkan lima status
    kelayakan dan menentukan apakah produk dapat dipublikasikan.

    \item \textbf{Food Score Decay} : Skor kelayakan (0--100) yang meluruh mengikuti sisa waktu
    rescue menuju \textit{expired}, ditampilkan \textit{real-time} beserta \textit{rescue timer}.

    \item \textbf{Keyword Classification dari Ulasan} : Klasifikasi keyword ulasan ke level
    Aman/Warning/Gawat dan akumulasi menjadi badge keamanan per UMKM.

    \item \textbf{Marketplace Customer} : Konsumen menelusuri produk rescue aktif, melihat Food
    Score, harga rescue, persentase diskon, dan badge keyword, serta menyaring dan mengurutkan produk.

    \item \textbf{Checkout Cashless via Midtrans} : Pembayaran \textit{cashless only} melalui Midtrans
    \textit{sandbox} dengan \textit{service fee} 5\% otomatis, pembuatan \textit{pickup code} setelah
    status \textit{Paid}, dan reservasi stok berbatas waktu.

    \item \textbf{Manajemen Pesanan dan Pickup} : Pencatatan pesanan dan siklus status (Created,
    Payment Pending, Paid, Ready for Pickup, Completed, No-Show, dll.) serta validasi \textit{pickup
    code} oleh UMKM.

    \item \textbf{Dashboard Analitik \& Insight UMKM} : Ringkasan penjualan, produk terlaris, tren,
    metrik listing (views, CTR, \textit{conversion rate}), serta \textit{insight} rating dan keyword.

    \item \textbf{Slot Iklan} : UMKM memasang iklan produk (paket berdurasi tetap) dan pihak ketiga
    dapat beriklan; seluruhnya melalui \textit{approval} admin.

    \item \textbf{Ulasan dan Rating} : Konsumen memberi rating dan ulasan dengan input keyword
    setelah transaksi selesai.

    \item \textbf{Help Center} : Customer melaporkan kendala (produk tidak tersedia/tidak sesuai,
    UMKM tidak merespons, masalah pembayaran) untuk ditangani admin.

    \item \textbf{Waste Log} : UMKM mencatat makanan tidak layak/tidak terjual sebagai bahan
    evaluasi produksi.

    \item \textbf{Konsol Admin} : Verifikasi UMKM, verifikasi Mitra Donasi (\textit{approve/reject}),
    moderasi listing dan iklan, dashboard keuangan (\textit{service fee} + iklan), dan \textit{export}
    laporan (CSV/Excel/PDF).

    \item \textbf{Mitra Donasi} : Registrasi organisasi non-profit dan verifikasi oleh admin untuk
    menerima makanan surplus.
\end{enumerate}

\subsection{Status Implementasi Fitur MVP}

Tabel~\ref{tab:mvp_scope} merangkum ruang lingkup MVP beserta status implementasi saat proposal ini
disusun. Fitur berstatus \emph{Sudah berjalan} telah tersedia pada basis kode (\textit{frontend}
Next.js dan/atau \textit{backend} Go/Fiber), sedangkan \emph{Direncanakan} merupakan bagian MVP yang
sedang/akan diselesaikan dalam \textit{sprint} berikutnya.

\begin{table}[htbp]
\centering
\caption{Ringkasan Ruang Lingkup dan Status MVP}
\label{tab:mvp_scope}
\small
\begin{tabularx}{\textwidth}{|>{\RaggedRight\arraybackslash}p{0.32\textwidth}|l|>{\RaggedRight\arraybackslash}X|}
\hline
\textbf{Fitur MVP} & \textbf{Status} & \textbf{Keterangan} \\
\hline
Manajemen produk UMKM & Sudah berjalan & CRUD produk reguler \& \textit{mystery box} \\
\hline
Food Score Decay & Sudah berjalan & Rule-based, meluruh terhadap waktu (\textit{frontend}) \\
\hline
Keyword classification ulasan & Sudah berjalan & Level Aman/Warning/Gawat (\textit{frontend}) \\
\hline
Marketplace customer & Sudah berjalan & Browse, filter, sort, Food Score live \\
\hline
Manajemen pesanan & Sudah berjalan & Pencatatan \& update status \\
\hline
Dashboard analitik \& insight UMKM & Sudah berjalan & Penjualan, tren, insight, metrik listing \\
\hline
Slot iklan UMKM & Sudah berjalan & Paket iklan berdurasi tetap \\
\hline
Ulasan \& rating & Sudah berjalan & Dengan input keyword \\
\hline
Autentikasi \& role (4 peran) & Direncanakan & JWT + \textit{password hashing}, RBAC \\
\hline
Food Trust Index (input UMKM) & Direncanakan & Form penilaian $\rightarrow$ 5 status kelayakan \\
\hline
Checkout cashless Midtrans + \textit{service fee} 5\% & Direncanakan & Midtrans \textit{sandbox}, \textit{webhook}, \textit{pickup code} \\
\hline
Help Center & Direncanakan & Tiket bantuan customer $\rightarrow$ admin \\
\hline
Waste Log & Direncanakan & Pencatatan makanan tidak terjual \\
\hline
Konsol admin \& keuangan & Direncanakan & Verifikasi, moderasi, revenue, \textit{export} \\
\hline
Mitra Donasi & Direncanakan & Registrasi + verifikasi admin \\
\hline
\end{tabularx}
\end{table}

\subsection{Di Luar Ruang Lingkup MVP (Out-of-Scope)}

Kapabilitas berikut secara sadar dikeluarkan dari MVP dan diposisikan sebagai \textit{roadmap}
pengembangan lanjutan setelah MVP tervalidasi:

\begin{itemize}[leftmargin=*]
    \item Mode \textit{production} Midtrans dan perluasan kanal pembayaran (MVP menggunakan
    \textit{sandbox} saja; tidak ada pembayaran tunai/COD).
    \item Layanan \textit{delivery}/pengantaran (MVP hanya \textit{self-pickup}).
    \item \textit{Dynamic pricing} otomatis dan prediksi permintaan berbasis \textit{machine learning}
    (MVP menyediakan rekomendasi diskon \textit{rule-based}, harga akhir ditetapkan UMKM).
    \item Pengelolaan limbah eksternal (MVP membatasi pada pencatatan internal via \textit{Waste Log}).
    \item Komunikasi \textit{real-time} (\textit{push notification}, \textit{websocket}) dan aplikasi
    \textit{native mobile}.
    \item Voucher/promo dan integrasi POS UMKM.
\end{itemize}

\subsection{Asumsi dan Batasan}

\begin{itemize}[leftmargin=*]
    \item Aplikasi berbasis web \textit{responsive}/PWA-ready dan diakses melalui peramban modern
    (\textit{mobile-first}), bukan aplikasi \textit{native}.
    \item Food Trust Index dan Food Score merupakan indikator kelayakan berbasis aturan dan input
    UMKM, bukan hasil uji laboratorium; keputusan akhir konsumsi tetap pada pengguna.
    \item Integrasi Midtrans menggunakan mode \textit{sandbox} untuk demo dan pengujian, bukan
    transaksi uang asli.
    \item Estimasi dampak lingkungan (makanan terselamatkan, emisi) bersifat estimatif untuk edukasi,
    bukan klaim terverifikasi pihak ketiga.
    \item Data UMKM, customer, produk, dan transaksi dapat menggunakan \textit{dummy data} realistis
    untuk keperluan demonstrasi dan pengujian.
    \item Role Mitra Donasi pada MVP terbatas pada registrasi dan verifikasi oleh admin.
\end{itemize}
